% Generated draft: requires mathematical and source review.
\documentclass[11pt]{article}
\usepackage[margin=25mm]{geometry}
\usepackage{fontspec}
\setmainfont{Latin Modern Roman}
\usepackage{amsmath,amssymb,mathtools,unicode-math}
\setmathfont{Latin Modern Math}
\usepackage{enumitem,needspace,xurl,hyperref}
\setlength{\parindent}{0pt}
\setlength{\parskip}{5pt}
\setlength{\emergencystretch}{3em}
\newcommand{\N}{\mathbb{N}}
\newcommand{\Z}{\mathbb{Z}}
\newcommand{\Q}{\mathbb{Q}}
\newcommand{\R}{\mathbb{R}}
\newcommand{\C}{\mathbb{C}}
\newcommand{\F}{\mathbb{F}}
\newcommand{\A}{\mathbb{A}}
\newcommand{\PP}{\mathbb{P}}
\newcommand{\E}{\mathbb{E}}
\newcommand{\eps}{\varepsilon}
\newcommand{\dd}{\,\mathrm{d}}
\begin{document}
\section*{UnsolvedMath ID: 20002089}
\textbf{Source problem code:} catalogue entry
\subsection{Definitions and mathematical statement}
Eliashberg's Question 1.31 asks to define a higher-genus version of relative contact homology as part of understanding the Gopakumar--Vafa picture through symplectic field theory. The recorder explicitly notes missing much of the discussion. The relative objects, coefficient ring, analytic foundations and comparison with Gopakumar--Vafa theory are not specified. ``Sixth day'' is the next heading, not part of the question.

One common intended setting would fix a contact manifold \((Y,\alpha)\), Legendrian boundary data \(\Lambda\), and adapted holomorphic curves in \(\R\times Y\), with boundary on \(\R\times\Lambda\), boundary punctures at Reeb chords and interior punctures at Reeb orbits, allowing higher genus. This is an explicit model rather than a uniquely recovered definition. Contact homology, genus-zero rational SFT, absolute full SFT, a relative open/closed theory, and closed Gopakumar--Vafa invariants must remain distinct.

The supplied aid reports a topological gluing ledger. Let the vertices of a gluing graph carry compactified oriented surface domains \(\Sigma_v\). Smooth an interior matching by removing two disks and inserting a neck, and a boundary matching by the corresponding boundary-interval gluing. If there are \(E_{\mathrm{int}}\) interior and \(E_{\partial}\) boundary contractions, then
\[
\chi(\Sigma_{\mathrm{out}})=\sum_v\chi(\Sigma_v)-2E_{\mathrm{int}}-E_{\partial},
\]
\[
g_s^{-\chi(\Sigma_{\mathrm{out}})}
=\left(\prod_vg_s^{-\chi(\Sigma_v)}\right)
g_s^{2E_{\mathrm{int}}+E_{\partial}}.
\]
Thus an interior contraction costs \(g_s^2\), while a boundary contraction costs \(g_s\). This counts the smoothed domain, not mere identification of two nodal points. Puncture/capping conventions must be fixed consistently before applying these powers to a Hamiltonian.

The aid's formal algebraic statement is conditional. Let \((\mathcal A,\star)\) be a complete graded associative algebra over characteristic zero with convergent formal operations, and let \(H\) be odd. Define the unscaled inner derivation \(D=\operatorname{ad}_\star H\), using the graded commutator. Then
\[
D^2=\operatorname{ad}_\star(H\star H).
\]
If the proposed geometric boundary equation is \(H\star H+B=0\), then
\[
D^2=-\operatorname{ad}_\star B.
\]
A noncentral residual boundary contribution therefore obstructs \(D^2=0\). A central residual may be invisible to the inner derivation, so \(D^2=0\) alone does not certify the full master equation. If a convention scales \(D\) by a coupling parameter, the corresponding squared prefactor must also be included.

To turn this ledger into the requested relative theory requires coherent oriented virtual chains, all codimension-one boundary strata, appropriate compactness and gluing, and a complete Hamiltonian compatible with the algebra. The aid expressly assumes rather than constructs these prerequisites. It supplies a termwise check on proposed theories, not a completed higher-genus invariant or a proof of the Gopakumar--Vafa correspondence.
\subsection{Short English statement}
Define a higher-genus relative contact/SFT theory and connect it to the Gopakumar--Vafa program. The supplied aid gives only a gluing-power ledger and a conditional master-equation obstruction. It does not provide the required virtual-chain construction or identify the different theories.
\subsection{Sources}
\begin{itemize}
\item Eliashberg's discussion, Question 1.31, AIM workshop \emph{Holomorphic curves in contact geometry}: \url{https://aimath.org/WWN/contactgeom2/contactgeom2.pdf}. The source explicitly includes the recorder's qualification ``I missed a lot here.''
\item The source names the Gopakumar--Vafa picture without a detailed reference. Local registry \texttt{data/unsolvedmath/problems.json}, permanent ID 20002089, code \texttt{catalogue entry}; canonical \texttt{aim-geometry-notes.json}, \texttt{notes[426]}; updated 14 August 2026. The original literature field is empty.
\item Attempt 1: conditional result, partially solved status, rigor label conditional, novelty confidence low. The machine-generated research aid is not peer reviewed and requires independent verification before citation.
\end{itemize}
\end{document}
